\section{Whole-partner support forces a second large alternating layer}
\label{sec:standard-partner-support-gap}

\paragraph{Author freeze and provenance.}
D-TD44, reserved question Q2ez, is a new derivation by
\texttt{/root/sixfolds} in the remaining nonfaithful-partner branch.
Root requested a bound using the entire original recursive partner,
instead of another coefficient-dimension refinement of HD-SS11.
The support-weighted Riemann--Hurwitz inequality and the two elementary
case divisions below are supplied in full.  The source, exact-field,
and natural-polarization interfaces are the accepted HD-SS01--11;
the complete factor-algebra interface is D-TD11--13 and HD-FT21.
Complete author written-proof, source-interface, checker and first-output
readings pass.  Independent and root whole-proof reviews also pass.
No source or rank-six budget for the partner alone is assumed.

\paragraph{Actual source and support parameter.}
Let \(C\to\mathbf P^1\) be an actual connected finite J-Galois cover.
Suppose the faithful absolutely irreducible coefficient representation
has an actual finite linear factorization
\[
 V=(\operatorname{Std}_n\circ x)\otimes U,\qquad
 x:J\twoheadrightarrow P\in\{A_n,S_n\},\quad n\ge5,
 \qquad b=\dim U\ge2,\qquad \mu_V=6.
 \tag{SPG.1}
\]
Assume V has a model over its exact imaginary quadratic character field K.
For the first assertion, suppose \(K_0=\ker U\) is transitive on the
natural letters and \(\mu=\mu_U>0\).  Use the actual quotients
\[
 \begin{gathered}
 H=x^{-1}(P_1),\quad N=H\cap K_0,\quad
 X=C/H,\quad Y=C/N,\quad C_U=C/K_0,\quad Q=U(J),\\
 h=g(X),\qquad H/N\simeq Q,\qquad \tau=U|_H.
 \end{gathered}
 \tag{SPG.2}
\]
Here \(P_1\) is a natural letter stabilizer.  The subgroup \(K_0\)
and coefficient field K are different objects.

Let \(M\ge80\) be a support denominator for the entire U-image,
in the precise sense that every nonscalar operator S in that image satisfies
\[
 \nu_U(S):=b-\max_{z\in\mathbf C^*}\dim\ker(S-zI)
                  \ \ge\ \frac bM.
 \tag{SPG.3}
\]
The value of M is not a dimension bound for U or for any of its pieces.

\begin{theorem}[D-TD44]
Under (SPG.1)--(SPG.3), including transitivity and \(\mu>0\),
\[
                              n\le108M-12.
 \tag{SPG.4}
\]
This assertion needs neither period rank nine nor minimality.

Now suppose the original coefficient image is primitive, and fix one
complete D-TD12 decomposition in which the displayed standard factor is
a separately stable singleton extraction layer.  Remove just that layer.
Let \(m_{*,U}\) be the largest natural degree of an alternating simple
quotient among all quasisimple factors in the entire remaining original
decomposition, with value one if there are none.  Then (SPG.3) holds with
\[
                         M_U=\max\{80,m_{*,U}-1\}.
 \tag{SPG.5}
\]
There is no omitted multiplicity space and no bound on the dimensions
or number of the remaining nonalternating or Heisenberg factors.

Finally suppose \(\ker U\ne1\), the sixfold variation has signature
\((3,3)\) and actual period rank nine, and the original primitive factor
carries its restricted Jacobian polarization.  Suppose also that the
source has the smallest coefficient degree among actual exact-K,
rank-nine sources of this natural signed polarization determinant type,
as in HD-SS06.  Then
\[
 n\le108\max\{80,m_{*,U}-1\}-12,
 \qquad
 n>8628\ \Longrightarrow\quad
       m_{*,U}\ge\left\lceil\frac{n+120}{108}\right\rceil.
 \tag{SPG.6}
\]
The factor detected by the last inequality is an actual remaining
quasisimple component of the fixed original decomposition.  It is not
asserted to be standard, separately stable, or an ordinary rather than
spin constituent; it may belong to a component-permuting layer.
\end{theorem}

\begin{proof}
\paragraph{Step 1: the actual lower quotient and the whole support floor.}
Transitivity gives \(J=K_0H\), \([K_0:N]=n\), and
\(H/N\simeq J/K_0\simeq Q\).  Thus \(Y\to X\) is an actual
connected Q-Galois cover, and \(Y\to C_U\) is an actual degree-n map.
It need not be Galois.  HD-SS01 proves that U and \(\tau\) are irreducible.
The stronger transitive-kernel part of HD-SS06 constructs their models
on the same exact field K; it does not substitute a larger splitting field.
In particular \(\tau\), regarded as a representation of Q, is faithful
and has the same matrices as the entire U-image.

Every genuine inertia operator of \(\tau\) is therefore nonidentity.
If it is nonscalar, its actual moving rank at target one is at least
its projective support in (SPG.3).  If it is a nonidentity scalar,
its moving rank is b and satisfies the same inequality.  Consequently,
writing o for its actual inertia order and r for its moving rank,
\[
             r\ge\frac bM,\qquad
             1-\frac1o\le1\le\frac Mb\,r.
 \tag{SPG.7}
\]
This uses every descended inertia, not just an original generator.
Every such matrix is a power of an original matrix, still in Q.
Neither an original lift order nor a projective order replaces o.

Positive \(\mu\) implies \(g(C_U)\ge2\): genus zero has no
cohomology, while in genus one the holomorphic and antiholomorphic
lines are invariant, so no irreducible constituent has degree \(b\ge2\).
Let s be the number of genuine branch points of the actual Q-cover
\(C_U\to\mathbf P^1\), with sorted orders \(d_1,\ldots,d_s\).
Its positive orbifold excess and the exact nontrivial-U Euler identity are
\[
 \Delta=-2+\sum_{i=1}^s(1-1/d_i)>0,\qquad
 D'=6+\mu-2b(h-1),\qquad
 \mu=(s-2)b-\sum_{i=1}^s f_i,
 \tag{SPG.8}
\]
where \(f_i\) are the U-fixed-space dimensions at those original
Q-inertias, and \(D'\) is the total moving rank on \(Y\to X\).
These are the complete identities of HD-SS02 and HD-SS08; the extra
cohomology \(\mu\) is retained.

Riemann--Hurwitz for \(Y\to C_U\), and then for \(Y\to X\),
together with (SPG.7), gives
\[
 \begin{split}
 n\Delta&\le\frac{2g(Y)-2}{|Q|}
       =2h-2+\sum_j(1-1/o_j)\\
       &\le2h-2+\frac Mb\bigl(6+\mu-2b(h-1)\bigr)\\
       &\le-2+\frac Mb(6+\mu+2b).
 \end{split}
 \tag{SPG.9}
\]
The last step is valid because \(2-2M<0\) and \(h\ge0\).
All maps are quotients or maps in the same original cover.  No
sphere-base assertion is applied to an arbitrary subgroup or constituent.

\paragraph{Step 2: every nonexceptional Q-signature.}
We reproduce the exact rows needed from the complete HD-SS08 table.
A hyperbolic three-mark signature with smallest order two has
\(\Delta\ge1/24\), unless it is \((2,3,7)\); also
\(\mu\le b-1\), since its involution is nonscalar and has a
nonzero fixed space.  If the smallest order is at least three,
\(\Delta\ge1/12\) and \(\mu\le b\).  Four genuine marks
have \(\Delta\ge1/6\) and \(\mu\le2b\).  For \(s\ge5\),
\(\Delta\ge(s-4)/2\) and \(\mu\le(s-2)b\).
The sorted-order proof and the nonscalar-involution proof are exactly
those of HD-SS08; thus these rows exhaust all positive signatures
other than \((2,3,7)\).

Substituting these values into (SPG.9) gives the following new bounds.
For comparison the middle column retains the old all-genus bounds of
HD-SS08, which use the integral rank bound one rather than \(b/M\).
\[
\begin{array}{c|c|c|c}
\text{actual signature case}
 &\text{old bound on }n&\text{new bound on }n
 &\text{combined bound on }n\\ \hline
s=3,\ d_1=2,\ (d_i)\ne(2,3,7)
 &72b+72&72M+120M/b-48&72M+72\\
s=3,\ d_1\ge3
 &36b+48&36M+72M/b-24&36M+48\\
s=4
 &24b+24&24M+36M/b-12&24M+24\\
s\ge5
 &10b+8&10M+12M/b-4&10M+8
\end{array}
 \tag{SPG.10}
\]
Here is the last row, including its uniformity in s:
\[
 n\le\frac{2Ms+12M/b-4}{s-4}
   =2M+\frac{8M+12M/b-4}{s-4}
   \le10M+12M/b-4.
 \tag{SPG.11}
\]
The numerator in the last fraction is positive, so its maximum is
at \(s=5\).  For every row of (SPG.10), if \(b\le M\) use
the old bound; if \(b\ge M\) use the new bound.  At \(b=M\)
they agree and yield exactly the last column.  This is valid for
all \(b\ge2\), without an asymptotic or large-dimension assumption.

\paragraph{Step 3: the exceptional triangle, with actual inertia orders.}
Suppose the Q-signature is \((2,3,7)\).  The exact quadratic-model
calculation HD-SS09 proves \(\mu\le b-\varepsilon_K(b)\le b-2\).
Since \(\mu>0\), it also gives \(b\ge3\).
HD-SS11 already proves the all-genus inequality \(n\le54b-12\).
We now combine it with the support denominator.

The connected cover \(Y\to X\) is the normalized base change
of \(C_U\to\mathbf P^1\).  More explicitly, at an original
inertia generator g with U-image q of order d, a cycle of length e
in its natural permutation gives local operator conjugate to \(q^e\)
and order \(d/\gcd(d,e)\).  Identity q produces no genuine mark.
Since \(d\in\{2,3,7\}\) and these orders are prime, every
genuine inertia on Y has order two, three, or seven.  This is
HD-SS11's local normalization calculation, including disappearing
original marks, and does not use the orders of the J-lifts.
Together with (SPG.7) it improves the weight bound to
\[
          1-\frac1o\le\frac67\le\frac{6M}{7b}\,r,
 \qquad o\in\{2,3,7\},\quad r\ge b/M.
 \tag{SPG.12}
\]
Thus the same actual Riemann--Hurwitz calculation gives
\[
 \begin{split}
 \frac n{42}
 &\le2h-2+\frac{6M}{7b}\bigl(6+\mu-2b(h-1)\bigr),\\
 n&\le72M+36M\mu/b+216M/b-84+(84-72M)h\\
  &\le108M+144M/b-84.
 \end{split}
 \tag{SPG.13}
\]
The second inequality uses \(\mu\le b-2\), \(h\ge0\), and
\(84-72M<0\).  No assumption that X is rational has been made.

If \(b\le2M\), HD-SS11 gives \(n\le54b-12\le108M-12\).
If \(b\ge2M\), (SPG.13) gives
\(n\le108M+72-84=108M-12\).  The two bounds agree at
\(b=2M\), so the triangle is covered for every b.  Each bound
in the last column of (SPG.10) is also at most \(108M-12\):
the four differences are \(36M-84\), \(72M-60\),
\(84M-36\), and \(98M-20\), all positive for \(M\ge80\).
This proves (SPG.4).

\paragraph{Step 4: M comes from the complete original complement.}
Fix the original D-TD12 decomposition now.  Its coarse extraction
algebras are individually stable; the selected singleton standard
layer is therefore a stable full matrix algebra.  Its commutant
is the entire U-algebra.  Removing that layer leaves a complete
tensor factorization of U by exactly the remaining quasisimple and
Heisenberg pieces, including every later multiplicity layer.
The original group permutes their embedded matrix algebras, and
an element fixing one algebra normalizes its indicated represented core.
These are the conclusions of D-TD12, not a new hypothesis that the
first generalized Fitting subgroup acts irreducibly.

This factor-algebra structure descends to the actual projective
image of U.  Indeed an original element scalar on U has trivial
conjugation on each embedded remaining matrix algebra.  Scalar
ambiguities in its factor representatives therefore do not alter
those algebra actions or their normalization.  Conversely, if
all remaining local operators are scalar and their factor algebras
are fixed, their full tensor product is scalar; there is no
additional multiplicity space outside this product.
The complete subcollection interface in D-TD17 is precisely this
original-image assertion.  It does not select a quotient map to
one component from a component-permuting layer.

The partner image is finite and primitive by D-TD11: a proper
invariant subspace or permuted direct-sum system for U would tensor
with the entire standard factor to give one for V.  More directly,
the remaining factor system already satisfies all the hypotheses
of HD-FT21 with \(A=2\) and with no arbitrary bounded-degree pieces.
Its local floors are \(1/80\) for nonalternating quasisimple cores,
\(1/(m-1)\) for alternating cores of degree \(m\ge7\),
\(1/46\) for degrees five and six, and at least \(1/4\)
for Heisenberg cores.  A nonscalar operator fixing every factor
algebra has at least the smallest local floor by spectral convolution.
If it moves a factor cycle, the accepted weighted tensor-cycle proof
gives the quarter floor, independently of all twists.  Consequently
these alternatives give (SPG.3) with exactly \(M=M_U\) in (SPG.5).
The maximum includes every remaining original layer.  No new recursion
with an uncontrolled maximal degree has been chosen, and there is
no assumption that the remaining components are individually stable.

\paragraph{Step 5: minimality, exact field and natural polarization.}
Suppose \(\ker U\ne1\).  Since V is faithful, the map
\(x|_{K_0}\) is injective: an element trivial on both
factors is trivial on V.  The nontrivial normal subgroup
\(x(K_0)\lhd P\) contains \(A_n\), so \(K_0\) is transitive.
This is the actual kernel dichotomy of HD-SS07 and uses no Jordan
bound.  It gives the exact same-field models and image equality in
Step 1, in particular \(U(H)=U(J)=Q\).

If \(\mu=0\), the post-vanishing portion of HD-SS06 applies with
its exact hypotheses and no size condition.  The actual multiplicity
variation for \(Y\to X\) is the direct sum of the U- and V-variations;
the first is zero.  Rank nine forces \(X\simeq\mathbf P^1\),
and the faithful irreducible degree-b Q-coefficient source has the
same exact K-model, signature \((3,3)\), and actual rank-nine variation.
The fixed multiplicity-one Rosati coefficient-line correspondence
and quotient norm/pullback identify its primitive sixfold with the
original one by a rational K-isogeny whose polarization scale
\(c_0\) is positive rational.  Its determinant scale is a norm because
\[
                           c_0^6
                   =\operatorname{Nm}_{K/\mathbf Q}(c_0^3).
 \tag{SPG.14}
\]
This is the actual restricted Jacobian polarization comparison of
(SSK.10)--(SSK.13), not a free choice of polarization or an assertion
about every isogeny.  The new coefficient degree \(b<(n-1)b\)
contradicts minimality in the stated comparison class.  Hence
\(\mu>0\), and Steps 1--4 give the first part of (SPG.6).

Finally \(108\cdot80-12=8628\).  For \(n>8628\), the maximum
in (SPG.5) must be supplied by \(m_{*,U}-1\).  Therefore
\[
                      n\le108m_{*,U}-120,
       \qquad m_{*,U}\ge\left\lceil\frac{n+120}{108}\right\rceil.
 \tag{SPG.15}
\]
The maximum is attained in a finite nonempty collection of actual
remaining factors.  This proves its existence with the stated scope
and completes the theorem.
\end{proof}

\paragraph{Strict remaining-gate reduction and limitations.}
For any fixed L, the condition that all remaining alternating simple
quotients have natural degree at most L now bounds the distinguished n
by \(108\max\{80,L-1\}-12\), independently of b.  In particular,
if no remaining alternating quotient has degree above81, then \(n\le8628\),
even if the partner contains arbitrarily many or arbitrarily large
nonalternating and Heisenberg pieces.  The earlier inequality
\(n\le72(b+1)\) permits arbitrarily large n in this parameter range
by taking b large; thus the new support constraint addresses a distinct
previously unbounded remaining gate.  No actual source satisfying an
excluded parameter configuration is asserted to exist.

The new factor need not be standard.  HD-MS03 applies only if it is
another separately stable ordinary standard layer; in that case it
also satisfies \(m\le(n+1)/2\).  No such upper bound is imposed here
on a nonstandard, spin, or component-permuting factor.  The proof does
not apply the sixfold moving-rank budget to U alone, discard a positive
U-cohomology summand, iterate a rank-nine descent with that summand
present, or give a global upper bound on n.  It preserves the exact
coefficient field, source family, and natural signed determinant only
in the explicit minimality step above.  No unrefereed theorem or
additional generic-endomorphism assumption is introduced.
